% §4.24 -- Iter 203 -- Cross-prompt pooling stress test (not empirical G')
\label{sec:p7-iter203-empirical-gprime}

\subsection{Cross-Prompt Pooling Is Not a Larger-Group Counterfactual}

An earlier revision described this analysis as an ``exact empirical $G'$''
counterfactual. That description was incorrect. The N2 tensor contains 16
different prompts with eight rollouts each. Concatenating two prompt rows creates
a 16-reward pool, but GRPO group-relative advantages are defined among rollouts
from the \emph{same prompt}. Cross-prompt pooling changes the estimand: it mixes
prompt difficulty and can manufacture contrast even when both original groups
are individually uniform.

For auditability we retain the descriptive pooling statistic
\[
\widehat{\mathrm{GU}}_{\mathrm{pool}}(G')=
\frac{1}{N_{G'}}\sum_{j=1}^{N_{G'}}
\mathbf{1}\{0<k_j^{(G')}<G'\},
\]
where contiguous prompt rows are concatenated to form pools of size
$G'\in\{16,32,64\}$. On the logged tensor its pooled means rise from $0.271$
at the factual within-prompt $G=8$ statistic to $0.547$, $0.784$, and $0.938$
after cross-prompt pooling. This monotonic rise is expected from heterogeneity:
combining an easy prompt with a hard prompt produces a mixed pool even if neither
prompt would produce within-prompt contrast at larger $G$.

\paragraph{What the result supports.}
The stress test demonstrates that a batch-level contrast statistic is highly
sensitive to whether prompt identity is preserved. It is useful as a negative
control for implementations and analyses that accidentally aggregate rewards
across prompts.

\paragraph{What the result does not support.}
It does not estimate $\mathrm{GU}(16)$ for any prompt, does not prescribe
escalating $G=8$ to $G=16$, and provides no controller-performance evidence.
A valid empirical $G'$ test requires fresh within-prompt rollouts under the same
policy, with total rollout tokens matched across controller and static arms.

\paragraph{Connection to the boundary audit.}
The valid 2,560-observation controller trace should guide the next experiment:
1,723 of 1,867 fires (92.3\%) occur on all-correct groups and 144 on all-wrong
groups. An all-wrong-only policy would remove 92.3\% of escalations, but whether
it preserves held-out performance is unknown. The preregistered comparison must
therefore include symmetric escalation, all-wrong-only retry, solved retirement
or distillation, and static $G=16$ under a common token budget.
